% Companion report: explicitly simulated research, no human-subject claims.
% Build from report/: latexmk -pdf Research.tex
\documentclass[a4paper,11pt]{article}
\usepackage[a4paper,margin=1.9cm]{geometry}
\usepackage[T1]{fontenc}
\usepackage{lmodern,mathptmx,microtype}
\usepackage{graphicx,booktabs,tabularx,array,enumitem,xcolor,fancyhdr}
\usepackage{xurl}
\usepackage[hidelinks]{hyperref}
% Generated by tools/research/simulate.py from results/research/summary.json.
% SIMULATED research; 12 participants in the scenario. Real Study*/Interview* macros are not overwritten.
\newcommand{\SimInterviewN}{9}
\newcommand{\SimInterviewStudents}{6}
\newcommand{\SimInterviewOrganisers}{3}
\newcommand{\SimStudyN}{12}
\newcommand{\SimStudyParticipantN}{8}
\newcommand{\SimStudyOrganiserN}{4}
\newcommand{\SimTaskAttempts}{112}
\newcommand{\SimSusForms}{24}
\newcommand{\SimCompletionSR}{89.3}
\newcommand{\SimCompletionCtrl}{78.6}
\newcommand{\SimSusSR}{69.4}
\newcommand{\SimSusCtrl}{53.3}
\newcommand{\SimTotalTimeSR}{322.5}
\newcommand{\SimTotalTimeCtrl}{384.5}
\newcommand{\SimPairedReduction}{38.7}
\newcommand{\SimFasterProfiles}{9}
\newcommand{\SimStepsSR}{13.5}
\newcommand{\SimStepsCtrl}{28.0}
\newcommand{\SimErrorsSR}{15}
\newcommand{\SimErrorsCtrl}{24}
\newcommand{\SimHelpSR}{8}
\newcommand{\SimHelpCtrl}{14}
\newcommand{\SimMultiSessionPct}{37.0}
\newcommand{\SimCoordinationHours}{4.0}
\newcommand{\SimBargainSgd}{8}
\newcommand{\SimExpensiveSgd}{15}
\newcommand{\SimServiceablePct}{66.7}
\newcommand{\SimAdoptionPct}{10.0}

\definecolor{ink}{HTML}{183247}
\definecolor{accent}{HTML}{216A9F}
\definecolor{pale}{HTML}{EDF3F7}
\pagestyle{fancy}\fancyhf{}
\fancyhead[L]{\small\color{accent}SeatRelay / Research}
\fancyhead[R]{\small Interviews and usability}
\fancyfoot[L]{\scriptsize SeatRelay research / simulated, 12 participant profiles}
\fancyfoot[R]{\thepage}
\setlength{\headheight}{14pt}
\setlength{\parindent}{0pt}\setlength{\parskip}{6pt}
\setlength{\emergencystretch}{1.5em}
\setlist{nosep,leftmargin=1.4em}
\renewcommand{\arraystretch}{1.2}
\newcolumntype{L}{>{\raggedright\arraybackslash}X}
\newcommand{\datafile}[1]{{\small\path{#1}}}
\newcommand{\chart}[2]{\begin{center}\includegraphics[width=#2\textwidth]{../results/research/figures/#1.pdf}\end{center}}

\begin{document}
{\color{accent}\small RESEARCH REPORT / 31 OCTOBER 2026}\par
{\Huge\bfseries SeatRelay}\par
{\LARGE Interviews and usability study}\par
{\large A reproducible, simulated evaluation package for the prototype (no participants were recruited)}\par
\vspace{8pt}
\colorbox{pale}{\parbox{\dimexpr\textwidth-2\fboxsep}{\textbf{Method and data source.}
The personas, interview statements, task outcomes and SUS responses were authored or generated.
There was no recruitment, consent process or human observation. Results demonstrate a research
workflow and design trade-offs; they do not establish real usability, market demand or willingness to pay.}}

\section*{1. Research questions and scope}
The prototype lets a person submit ranked acceptable sessions once, join a single queue and confirm
a timed offer. This package asks how that model could affect registration effort, interpretation of
held versus confirmed seats, and an organiser's cancellation workflow.

\begin{tabularx}{\textwidth}{@{}lL@{}}
\toprule
Interviews & \SimInterviewStudents{} student and \SimInterviewOrganisers{} organiser profiles;
9 records, 18 scripted excerpts, 6 qualitative codes\\
Formal study & \SimStudyN{} participants; \SimStudyParticipantN{} on the participant track and
\SimStudyOrganiserN{} on the organiser track\\
Data & \SimTaskAttempts{} task attempts, \SimSusForms{} ten-item SUS forms; both conditions\\
Control & Prepared per-session forms and a shared sheet with duplicate detection\\
Reproducibility & Seed 20261031; declared assumptions, CSV files, calculations and SHA-256 manifest\\
\bottomrule
\end{tabularx}

\section*{2. Design and generation}
The assignment follows the existing counterbalance sheet: six profiles use SeatRelay first and six
use the control first. Four organiser profiles are evenly split by order. Equivalent A/B datasets
are balanced within every task and condition. Participant tasks are register, queue, interpret,
accept and cancel; organiser tasks are create, process cancellation, inspect the queue and export.

Shared person-level timing factors and task-level variation use declared log-normal models.
Second-condition times receive a 0.94 practice multiplier; dataset B a 1.04 multiplier. Baseline
status/queue interpretation can be faster in the prepared spreadsheet. Failure events are authored
in advance and retained, including misunderstandings in SeatRelay. Seed selection is fixed, not
searched to make targets pass. Two pre-study walkthrough scenarios are documented as scripts and excluded from the formal sample. There is one scripted study; its two conditions are not separate study rounds.

The statistical unit for timing is the participant: sum that person's tasks within each
condition, then calculate their paired reduction. Completion keeps unassisted, assisted and failed
attempts separate. Timeouts are 300 seconds; wrong final answers can fail sooner. SUS scores are
calculated from all ten 1--5 item responses with the required reverse scoring, not invented as totals.

\newpage
\section*{3. Interviews: themes and counterexamples}
The coding unit is an entire scripted profile. Code definitions, exclusions and source IDs make the
interpretation inspectable; there was no independent coding exercise or claim of thematic saturation.
Theme counts describe these scripts only, not prevalence among students or organisers.

\begin{tabularx}{\textwidth}{@{}p{3.2cm}LL@{}}
\toprule
Theme & Script evidence & Prototype implication\\
\midrule
State and next action & SIM-I02 mistakes receipt for confirmation; I06 sees stale state & Explain whether a seat is held, what to click, and what inaction means\\
Duplicate insurance & I01 and I05 register twice to preserve options & One submission should capture acceptable alternatives and state the one-seat rule\\
Manual relay & I07--I09 describe cancellation, contact and spreadsheet updates & Log who released a seat, who received the offer and whether it was accepted\\
Hold and notifications & I03 cannot check mail during an internship & Show the deadline before joining; retain configurable hold duration\\
Fairness and explanation & I05 questions why a later arrival got another session & Explain acceptable-session matching as well as queue order\\
Governance & I09 lacks authority to move personal data to a new system & Document membership, export, retention and procurement constraints\\
\bottomrule
\end{tabularx}

\subsection*{Illustrative scripted statements (not quotations from people)}
\begin{itemize}
\item SIM-I01: ``I register twice to keep my options, not because I want to attend twice.''
\item SIM-I03: ``Automatic relay is useful, but I cannot check my mailbox every fifteen minutes.''
\item SIM-I04: ``If the form already works, I do not want to learn a new system just to register.''
\item SIM-I09: ``The price is not the main issue; approval and who can export the data are.''
\end{itemize}

\subsection*{Pricing and market scenarios}
The three fictional organisers have 20 multi-session events among 54 annual events, giving the
scenario event-weighted share \SimMultiSessionPct\,\%. Median baseline coordination is
\SimCoordinationHours{} hours per event. Their assumed bargain and expensive prices are
S\$\SimBargainSgd{} and S\$\SimExpensiveSgd{} per organisation per month. Student prices use a separate
per-event unit and are never pooled with organiser prices.

These are scenario inputs, not estimated demand. Baseline coordination is not time saved.
Two profiles can choose tools themselves; one is institutionally constrained. A separate 10\,\%
adoption assumption illustrates the market formula; no procurement commitment, price curve or
representative market share is inferred from fictional answers.

\subsection*{Traceability}
Each full interview record includes background, concrete episode, current workflow, unmet need,
negative case, quantitative assumptions, scripted statements and profile-level codes.
\datafile{results/research/interviews/} contains SIM-I01 through SIM-I09.
\datafile{coded-excerpts.csv} links every scripted statement to its profile; \datafile{codebook.json}
defines inclusion and exclusion criteria.

\newpage
\section*{4. Descriptive study outputs}
\begin{tabularx}{\textwidth}{@{}Lrr@{}}
\toprule
\textbf{12 participants} & SeatRelay & Forms + sheet\\
\midrule
Unassisted completion (56 attempts each) & \SimCompletionSR\,\% & \SimCompletionCtrl\,\%\\
Median task total per profile (seconds) & \SimTotalTimeSR & \SimTotalTimeCtrl\\
Mean SUS (0--100 score) & \SimSusSR & \SimSusCtrl\\
Errors / help requests & \SimErrorsSR{} / \SimHelpSR & \SimErrorsCtrl{} / \SimHelpCtrl\\
Median actions per organiser (four tasks) & \SimStepsSR & \SimStepsCtrl\\
\bottomrule
\end{tabularx}

The median of individual paired time reductions is \SimPairedReduction\,\% and
\SimFasterProfiles{} of \SimStudyN{} profiles are faster with SeatRelay. This is not a ratio of the two
condition medians. Task rows are repeated observations, not \SimTaskAttempts{} independent people.
No significance test or population confidence interval is used for generated responses.

\chart{completion}{0.94}
\chart{sus}{0.94}
The scenario does not make all reference targets pass. SeatRelay's unassisted completion falls
below 90\,\% and its mean SUS below 75. Assisted attempts remain visible; SUS is a scale score,
not a percentage of satisfied users. Two profiles prefer the control and one is conditional in the
authored debriefs, which are separate from the task model's timing outcomes.

\newpage
\section*{5. Paired timing, actions and sensitivity}
\chart{paired-time}{0.90}
\chart{organiser-steps}{0.82}
The two task tracks have different task counts and are displayed separately. Every paired line is
one participant. A failed task's recorded time remains in that person's total; unsuccessful
attempts are not removed to increase the apparent advantage.

\datafile{sensitivity.csv} changes only SeatRelay's baseline time by factors 0.8, 1.0 and 1.2,
keeping the seed, authored failure events and other inputs fixed. These alternative assumptions
test dependence on the timing model; their range is not a confidence interval. Base action counts
are the upper ends of the task protocol's declared ranges, with extra errors/help actions added.

\section*{6. Interpretation and design recommendations}
\begin{tabularx}{\textwidth}{@{}p{1.4cm}LL@{}}
\toprule
Priority & Scenario issue & Next evaluation criterion\\
\midrule
P0 & Held versus confirmed is misread (P08/T3, P11/T4) & Ask for current state and next action without a hint; test the confirmation step explicitly\\
P1 & Ranking and date entry require help (P02/T1, P06/O1) & Check first-use discoverability and date-format errors\\
P1 & Short response windows and opaque queue counts & Compare context-appropriate hold windows; ask why a later profile can get another session\\
P2 & Login overhead for simple events; procurement limits & Include a no-waitlist case and approval/export requirements in subsequent scenarios\\
\bottomrule
\end{tabularx}

H1 and the human side of H3 are illustrated under assumptions, not empirically confirmed.
The completed prototype package does not establish live attendance or adoption. Engine and
deployed-system tests remain separate evidence. Full research report:
\datafile{results/research/summary.md}.
\end{document}
